\scanpage{18}
% L1
foncteurs sur un topos $\widetilde C$ commutant aux $\varinjlim$

% L2–3
Soient $C$ un $\mathcal U$-site, $E$ une catégorie avec petites $\varinjlim$, considérons les catégories

% L4
\[
\underline{\mathrm{Hom}}'(\widetilde C,E)\qquad\underline{\mathrm{Hom}}\,(C,E)
\]

% L5–6
\deleted{et} \deleted{Soit} $\underline{\mathrm{Hom}}'(\widetilde C,E)$ \deleted{désignant} la sous-catégorie pleine de $\underline{\mathrm{Hom}}$ formée des foncteurs qui commutent aux petites $\varinjlim$. On va définir des foncteurs

% L7
\diagram{\begin{tikzcd}[ampersand replacement=\&,column sep=large]
\underline{\mathrm{Hom}}\,(\widetilde C,E)\arrow[r,shift left,"\alpha"] \& \underline{\mathrm{Hom}}\,(C,E)\arrow[l,shift left,"\beta"]
\end{tikzcd}}

% L8
en posant

% L9–10
\[
\begin{aligned}
\alpha(u)&=u\circ\varepsilon_C
 &&(\varepsilon_C\colon C\to\widetilde C\ \text{le foncteur canonique})\\
\beta(u_0)&=[F\to\mathop{\underline{\lim}}_{C/F}u_0(x)]
 &&(=\widehat\beta\mid\widetilde C,\ \text{où }\widehat\beta\colon\widehat C\longrightarrow E\ \text{est}\\
&&&\text{le prolongement can. de }\beta\ \text{défini}\\
&&&\text{par commutation aux }\varinjlim)
\end{aligned}
\]

% L11–14-right
\diagram{\begin{tikzcd}[ampersand replacement=\&,row sep=small,execute at end picture={\draw (current bounding box.north east) -- (current bounding box.north west) -- (current bounding box.south west) -- (current bounding box.south east);}]
C\arrow[r,hook,"{j_0}"'] \& \widehat C\arrow[dl,bend left,"a"] \& \varepsilon_C=aj_0\\
\widetilde C\arrow[ur,hook,"{i_0}"] \& \&
\end{tikzcd}}

% L11–12
Soit $u\in\mathrm{Ob}\,\underline{\mathrm{Hom}}\,(\unsure{\overline C|\widetilde C},E)$, on a un hom

% L13
\[
i_u\colon\beta\alpha(u)\longrightarrow u
\]

% L14
donné \unsure{pour|par} $\forall F\in\mathrm{Ob}\,\widetilde C$ par

% L15–16
\diagram{\begin{tikzcd}[ampersand replacement=\&,row sep=small]
{i_u(F)\colon} \& (\beta(\alpha(u)))(F)\arrow[r] \& u(F)\\
\& \varinjlim_{C/F}u(\varepsilon_C(x))\arrow[u,equal] \&
\end{tikzcd}}

% L16–18-right
déduit de l'isom can. $\mathop{\varinjlim}_{C/F\,(\widetilde C)}X\xrightarrow{\sim} F$ en appliquant $u$

% L19–20
On trouve que l'hom $i_u(F)$ est bien fonctoriel en $F$, d'où $i_u$, et que $i_u$ est fonctoriel en $u$, d'où

% L21
\[
\boxed{\beta\alpha\xrightarrow{i}\mathrm{id}_{\underline{\mathrm{Hom}}(\widetilde C,E)}}
\]

% L22–24
Notons tout de suite que \deleted{\unsure{$u$}} $i_u\colon\beta\alpha(u)\longrightarrow u$ est un isom \unsure{sss|ssi} $u\colon\widetilde C\longrightarrow E$ commute aux $\varinjlim$ du type $\varinjlim_{C/F}X\xrightarrow{\sim}F$ ($F\in\mathrm{Ob}\,\widetilde C$) ; Il \uline{suffit} \inserted{donc} que $u\in\mathrm{Ob}\,\underline{\mathrm{Hom}}'(\widetilde C,E)$.

% L25
Soit $u_0\in\underline{\mathrm{Hom}}\,(C,E)$, on va définir

% L26
\[
j_{u_0}\colon\deleted{\unsure{\alpha\beta}}u_0\deleted{\unsure{\}}}\unsure{\longleftrightarrow|\longrightarrow}\alpha\beta(u_0)
\]

% L27
\deleted{en posant}, \inserted{donné} pour $Z\in\mathrm{Ob}\,C$, par

% L28–30
\diagram{\begin{tikzcd}[ampersand replacement=\&,row sep=small]
{j_{\unsure{u_0|0}}(Z)\colon}\ \deleted{\unsure{(\alpha(\beta(u_0)))(Z)\longrightarrow}}
 \& u_0(Z)\arrow[r] \& (\alpha(\beta(u_0)))(Z)\\
\& \& \beta(u_0)(\varepsilon_C(Z))\arrow[u,equal]\\
\& \& \varinjlim_{C/\varepsilon_C(Z)}u_0(X)\arrow[u,equal]
\end{tikzcd}}

% L31–32
en notant que $(Z,\deleted{\unsure{\varepsilon_C(Z)}}\inserted{\text{morphisme can. de } Z\text{ dans }\varepsilon_C(Z)}\unsure{\}})$ est un objet de $C/\varepsilon_C(Z)$, soit $X_0$, et que par suite $u_0(X_0)=u_0(Z)$ s'en va dans la $\varinjlim$.

% foot
% archive sheet number: 17
